\documentclass[12pt]{article}
\usepackage[utf8]{inputenc}
\usepackage[T1]{fontenc}
\usepackage[margin=1in]{geometry}
\usepackage{setspace}
\usepackage{xcolor}
\usepackage[colorlinks=true, urlcolor=blue!60!black, linkcolor=blue!60!black]{hyperref}
\setlength{\parindent}{2em}
\setlength{\parskip}{1.25em}
\renewcommand{\baselinestretch}{1.0}
\begin{document}
\begin{center}
{\LARGE \textit{Fauntleroy v. Lum}, 210 U.S. 230 (1908)\par}
\vspace{0.8em}
{\large Case Brief\par}
\end{center}
\vspace{1em}

\section{Memory Jogger}
A Missouri judgment on an award arising from a Mississippi cotton-futures contract that was illegal in Mississippi had to be enforced in Mississippi: the Court held 5--4, per Justice Holmes, that full faith and credit bars collateral attack on a sister-state judgment for errors of law, leaving only want of jurisdiction as a defense (\href{https://supreme.justia.com/cases/federal/us/210/230/}{Justia}).

\section{Detailed Case Facts}
Two Mississippi residents entered into a speculative cotton-futures contract in Mississippi without any intent to deliver the underlying commodity or pay its actual price. Under Mississippi's Annotated Code of 1892 (\S\S~1120, 1121, 2117), dealing in cotton futures was a criminal misdemeanor, and such contracts were declared void and ``shall not be enforced by any court.'' When the defendant, Lum, refused to pay his losses on the transaction, the parties submitted their dispute to arbitration in Mississippi, omitting the question of the contract's illegality from the submission. The arbitrator issued an award against Lum.

\section{Procedural History}
As Justice White's dissent recounts, the claimant first sued on the arbitration award in a Mississippi state court but voluntarily dismissed the action once the court's attention was drawn to the illegal and criminal character of the underlying futures contract. Later, catching Lum temporarily present in Missouri, the claimant served him personally and sued on the arbitration award in a Missouri state court. The Missouri trial judge refused to admit Lum's evidence that the underlying transaction was illegal and void under Mississippi law, instead instructing the jury to return a verdict for the plaintiff if the arbitration submission and unpaid award were proved. The jury found for the plaintiff, and the Missouri court entered final judgment.

The plaintiff assigned the Missouri judgment to Fauntleroy (his attorney), who then sued to enforce the Missouri judgment in a Mississippi state court. Lum pleaded the illegality of the original cotton-futures transaction and Mississippi's statutory bar. Fauntleroy demurred and filed replications invoking the Full Faith and Credit Clause (U.S. Const. art. IV, \S~1). The Supreme Court of Mississippi upheld Lum's defense and refused to enforce the Missouri judgment, concluding that Mississippi courts lacked authority to give effect to a judgment rooted in an illegal, criminal Mississippi gambling transaction. Fauntleroy sought review in the U.S. Supreme Court by writ of error.

\section{Judicial Votes}
The Supreme Court \textbf{reversed} the Supreme Court of Mississippi by a \textbf{5--4} vote. The case was argued on April 27--28, 1908, and decided on May 18, 1908 (\href{https://supreme.justia.com/cases/federal/us/210/230/}{Justia}):

\begin{itemize}
    \item \textbf{Majority Opinion:} Justice Holmes delivered the opinion of the Court, joined by Chief Justice Fuller and Justices Brewer, Peckham, and Moody.
    \item \textbf{Concurring Opinions:} None.
    \item \textbf{Dissenting Opinion:} Justice White filed a dissenting opinion, joined by Justices Harlan, McKenna, and Day.
\end{itemize}

\section{Holding}
\textbf{Questions Presented:}

\begin{enumerate}
    \item Does a state statute declaring that gambling contracts ``shall not be enforced by any court'' strip a state court of general jurisdiction of subject-matter jurisdiction to entertain an action on a sister-state judgment, or does it merely establish a substantive rule of decision?
    \item Under the Full Faith and Credit Clause (U.S. Const. art. IV, \S~1) and its implementing statute (1 Stat. 122), must a state court enforce a final money judgment rendered by a sister-state court with personal and subject-matter jurisdiction even if the rendering court misapplied the enforcing state's substantive law and enforced a claim contrary to the enforcing state's criminal statutes and public policy?
\end{enumerate}

\textbf{Holding:} \textbf{Yes} on the constitutional duty to enforce; \textbf{reversed} (5--4). The Supreme Court held that Mississippi's anti-futures statute merely laid down a substantive rule of decision rather than withdrawing subject-matter jurisdiction from Mississippi's courts of general jurisdiction. Because the Missouri court had undisputed personal and subject-matter jurisdiction, the Full Faith and Credit Clause required Mississippi to give the Missouri judgment the same conclusive effect it had in Missouri, even if the Missouri court misapprehended Mississippi law and enforced an obligation illegal under Mississippi's criminal laws and public policy.

\section{Analysis of Opinions}
\subsection{Majority Opinion (Justice Holmes)}

Writing for a five-Justice majority, Justice Holmes dismantled Mississippi's defense in two tightly structured steps: first separating jurisdiction from merits, and second affirming an ironclad rule of interstate judgment finality.

\subsection{Jurisdiction Versus Merits}
Lum argued that Mississippi's statute deprived Mississippi courts of jurisdiction, invoking \textit{Anglo-American Provision Co. v. Davis Provision Co.}, 191 U.S. 373 (1903), where New York lawfully closed its courts to suits between foreign corporations on out-of-state causes of action. Holmes rejected the analogy by drawing a sharp conceptual line between a court's \textit{power} (jurisdiction) and its \textit{duty} (merits). A court of general jurisdiction has power to decide a dispute even when substantive law forbids recovery; if it errs on the merits, the judgment is erroneous but not void. Construing ``shall not be enforced by any court'' as a jurisdictional ouster would make no more sense than treating the Statute of Frauds (``an action shall not be brought'') as jurisdictional. When a statute affects a court of general jurisdiction, courts must be ``slow to read ambiguous words'' as stripping power rather than fixing a rule of decision.

\subsection{Conclusiveness of Sister-State Judgments}
Turning to the merits of the Full Faith and Credit defense, Holmes anchored the Court's analysis in Chief Justice Marshall's classic formulation in \textit{Hampton v. M'Connel}, 16 U.S. (3 Wheat.) 234 (1818), and the Act of May 26, 1790 (1 Stat. 122): a state judgment must receive the same ``credit, validity, and effect'' in every sister state that it enjoys in the rendering state. Holmes brushed aside contrary language in \href{https://supreme.justia.com/cases/federal/us/127/265/}{\textit{Wisconsin v. Pelican Ins. Co.}, 127 U.S. 265 (1888)}, distinguishing it as a case about the Supreme Court's own Article III original jurisdiction over penal fines and dismissing its broader remarks about looking behind a judgment as mere dictum.

Once a rendering court has valid personal and subject-matter jurisdiction, its final judgment is conclusive as to all \textit{media concludendi} (the intermediate legal and factual steps leading to judgment) and cannot be collaterally impeached in another state for an error of law. Lum's proper remedy for Missouri's legal error was direct appeal within Missouri (and, if a federal question existed, to the U.S. Supreme Court), not collateral attack in Mississippi. Holmes dismissed fears of interstate humiliation, predicting that sister-state mistakes would be ``rare'' and noting that every state stands ``equally on both sides'' of mutual judgment enforcement.

\subsection{Dissenting Opinion (Justice White)}

White argued that the majority expanded the Full Faith and Credit Clause far beyond its founding-era purpose, warning that it would ``obliterate all state lines'' and empower any state to nullify the police powers and criminal statutes of its sisters. Relying on Joseph Story's \textit{Conflict of Laws} and precedents including \textit{Thompson v. Whitman}, 85 U.S. (18 Wall.) 457 (1874), \href{https://supreme.justia.com/cases/federal/us/127/265/}{\textit{Wisconsin v. Pelican Ins. Co.}, 127 U.S. 265 (1888)}, and \textit{Andrews v. Andrews}, 188 U.S. 14 (1903), White contended that the Clause merely made obligatory the preexisting principles of international comity, which never compelled a sovereign to enforce a foreign judgment validating a domestic criminal act.

White also challenged Holmes's reading of both \textit{Pelican} and \textit{Anglo-American Provision Co.} In White's view, \textit{Pelican} expressly authorized an enforcing court to look behind a judgment ``to ascertain whether the claim is really one of such a nature that the court is entitled to enforce it.'' Moreover, because the Mississippi Supreme Court had authoritatively construed its own anti-futures statute to deny complainants access to Mississippi courts, the U.S. Supreme Court was bound by that jurisdictional construction. Finally, White dismissed the intervening arbitration award as ``negligible,'' invoking the maxim that what is void for public policy cannot be cured by confirmation or acquiescence.

\subsection{Subsequent Developments}

\begin{itemize}

\item \textbf{Extension to Jurisdiction-Denying Statutes:} In \href{https://supreme.justia.com/cases/federal/us/252/411/}{\textit{Kenney v. Supreme Lodge of the World}, 252 U.S. 411 (1920)}, Holmes extended \textit{Fauntleroy} to hold that a state cannot evade Full Faith and Credit by the simple device of denying jurisdiction to its courts over out-of-state wrongful-death judgments.

\item \textbf{No Public-Policy Exception Reaffirmed:} In \href{https://supreme.justia.com/cases/federal/us/522/222/}{\textit{Baker v. General Motors Corp.}, 522 U.S. 222 (1998)}, Justice Ginsburg reaffirmed \textit{Fauntleroy}'s core rule that the Full Faith and Credit obligation for judgments is ``exacting'' and admits no public-policy exception, even while carving out a narrow limit: one state's equitable anti-suit or anti-testimony decree cannot control the evidentiary machinery of a sister state's courts in litigation involving nonparties.

\item \textbf{Applied to Adoption Decrees:} In \href{https://supreme.justia.com/cases/federal/us/577/464/}{\textit{V.L. v. E.L.}, 577 U.S. 404 (2016)} (per curiam), a unanimous Court summarily reversed the Alabama Supreme Court for refusing to recognize a Georgia second-parent adoption decree, directly quoting Holmes's rule in \textit{Fauntleroy} that courts must be ``slow to read ambiguous words'' in a sister state's statute as jurisdictional rather than merits-based.

\end{itemize}

\section{Examples: Future Applications}
\begin{enumerate}

\item \textbf{Same-Side (Judgment Must Be Enforced): An Illegal Sports Wager.} State A prohibits sports-betting debts as a misdemeanor and declares them ``unenforceable in any court.'' Resident X and Resident Y, both domiciled in State A, enter an illegal wager in State A. While Y attends a trade show in State B, X serves Y and obtains a final civil money judgment in State B's trial court, which mistakenly rejects Y's State A illegality defense. Y does not appeal in State B. When X seeks to enforce the judgment in State A, State A's courts refuse on public-policy grounds. Under \textit{Fauntleroy}, State A \textbf{must} enforce the judgment. State B had valid personal jurisdiction (in-state service) and subject-matter jurisdiction; its mistake of State A law had to be corrected on direct appeal in State B, not collaterally in State A.

\item \textbf{Same-Side (Judgment Must Be Enforced): Damages for Breach of a Non-Compete.} State C's statutes declare employee non-compete agreements void as against public policy. An employer sues a former employee in State D, where the employee worked and where non-competes are enforceable, and obtains a final money judgment for breach of a non-compete covering State C. The employee moves home to State C, and the employer seeks to enforce the judgment there. State C's courts may not refuse on the ground that the underlying agreement offends State C's policy. State D had jurisdiction over the parties and the claim, so under \textit{Fauntleroy} its judgment is conclusive even if it misjudged which State's law should govern; the employee's remedy was appeal in State D (\href{https://supreme.justia.com/cases/federal/us/210/230/}{Justia}).

\item \textbf{Opposite-Side (Enforcement May Be Refused): A State Civil Penalty Against an Absent Defendant.} Now alter the facts so that \textit{State B's Attorney General} obtains a \$100,000 statutory civil penalty in State B against Y (a State A resident who never entered State B and never appeared) for mailing regulated goods into State B, and seeks enforcement in State A. Here, State A has two potent \textit{Fauntleroy}-compatible defenses that flip the outcome: (1) \textit{lack of personal jurisdiction} in State B (which Holmes expressly preserved as a valid collateral attack, 210 U.S. at 237, and which Y did not waive by defaulting), and (2) the \textit{penal-judgment exception} (\href{https://supreme.justia.com/cases/federal/us/146/657/}{\textit{Huntington v. Attrill}, 146 U.S. 657 (1892)}), arguing that a state-brought punitive fine---unlike Fauntleroy's private contract/arbitration recovery---is a penal obligation that sister states need not execute.

\item \textbf{Opposite-Side (Decree Need Not Be Obeyed): An Injunction Against Testifying.} A third contrast arises under \href{https://supreme.justia.com/cases/federal/us/522/222/}{\textit{Baker v. General Motors Corp.}, 522 U.S. 222 (1998)}: if State B enters an injunction barring Y from testifying in State A on behalf of nonparty strangers, State A may disregard the order because one state's decree cannot hijack another state's evidentiary procedures against nonparties.

\item \textbf{Fence-Sitter (Uncertain): A Private Bounty Judgment.} State E authorizes any private person to sue anyone who assists a specified activity and to recover at least \$10,000 per violation, without proof of personal injury. A State E court, with personal jurisdiction over the defendant, enters such a judgment, and the plaintiff seeks to enforce it in State F, where the activity is lawful. \textit{Fauntleroy} forbids State F from refusing merely because the judgment offends its public policy. State F's best argument is the penal-judgment exception of \textit{Huntington v. Attrill}: the award looks like a fine for a public wrong. But the plaintiff is a private party collecting a civil award, which resembles the private recovery enforced in \textit{Fauntleroy}. Which side of that line a bounty judgment falls on is unsettled.

\end{enumerate}

\section{Critique}
\textbf{Forum-Shopping and Transient Jurisdiction:} Critics note that \textit{Fauntleroy} incentivizes strategic ``tag'' jurisdiction (serving a defendant passing through a favorable forum) to launder a claim barred in the parties' home state into an unimpeachable sister-state judgment.

\textbf{Asymmetry Between Laws and Judgments:} Modern doctrine creates a stark cliff: before judgment, Mississippi was free to apply its own anti-gambling statute; the instant Missouri entered judgment---even by glaring legal error---Mississippi's sovereign policy evaporated.

\textbf{Holmes's Optimism Tested:} Holmes predicted that deliberate sister-state overrides of another state's criminal or public-policy laws would be ``rare.'' Modern culture-war federalism (reproductive care, firearms liability, and non-compete bans) directly tests that assumption.

\section{Key Quotations}
All quotations are verified verbatim from \href{https://supreme.justia.com/cases/federal/us/210/230/}{Justia (210 U.S. 230)}:

\begin{quote}
``No doubt it sometimes may be difficult to decide whether certain words in a statute are directed to jurisdiction or to merits, but the distinction between the two is plain. One goes to the power, the other only to the duty, of the court.''
\end{quote}
\noindent --- Justice Holmes, 210 U.S. at 234--235 (\href{https://supreme.justia.com/cases/federal/us/210/230/}{Justia}).

\begin{quote}
``When it affects a court of general jurisdiction and deals with a matter upon which that court must pass, we naturally are slow to read ambiguous words as meaning to leave the judgment open to dispute or as intended to do more than to fix the rule by which the court should decide.''
\end{quote}
\noindent --- Justice Holmes, 210 U.S. at 235 (\href{https://supreme.justia.com/cases/federal/us/210/230/}{Justia}).

\begin{quote}
``A judgment is conclusive as to all the \textit{media concludendi}, \ldots{} and it needs no authority to show that it cannot be impeached either in or out of the state by showing that it was based upon a mistake of law. Of course, a want of jurisdiction over either the person or the subject matter might be shown. \ldots{} But, as the jurisdiction of the Missouri court is not open to dispute, the judgment cannot be impeached in Mississippi even if it went upon a misapprehension of the Mississippi law.''
\end{quote}
\noindent --- Justice Holmes, 210 U.S. at 237 (\href{https://supreme.justia.com/cases/federal/us/210/230/}{Justia}).

\begin{quote}
``Although not wishing in the slightest degree to weaken the operation of the due faith and credit clause as interpreted and applied from the beginning, it to me seems that this ruling so enlarges that clause as to cause it to obliterate all state lines, since the effect will be to endow each state with authority to overthrow the public policy and criminal statutes of the others, thereby depriving all of their lawful authority.''
\end{quote}
\noindent --- Justice White, dissenting, 210 U.S. at 239--240 (\href{https://supreme.justia.com/cases/federal/us/210/230/}{Justia}).


\end{document}
